\begin{table}[H]
\centering
\caption{Cross-model generalization and adaptation under disjoint calibration and shared held-out evaluation data.}
\label{tab:generalization}
\small
\begin{adjustbox}{max width=\textwidth}
\begin{tabular}{@{}llccccc@{}}
\toprule
Evaluation LLM & Router / Profile Source & Calibration Questions & F1 (\%) & Exact Match (\%) & Mean F1 Regret & Avg. Retrieval Calls \\
\midrule
\multirow{4}{*}{LLM 1}
& Original Adaptive-RAG router & \NA & \NA & \NA & \NA & \NA \\
& KNN calibrated on LLM 1 gains & \NA & \NA & \NA & \NA & \NA \\
& KNN calibrated on LLM 2 gains & \NA & \NA & \NA & \NA & \NA \\
& Model-matched distance-weighted KNN & \NA & \NA & \NA & \NA & \NA \\
\midrule
\multirow{4}{*}{LLM 2}
& Original Adaptive-RAG router & \NA & \NA & \NA & \NA & \NA \\
& KNN calibrated on LLM 1 gains & \NA & \NA & \NA & \NA & \NA \\
& KNN calibrated on LLM 2 gains & \NA & \NA & \NA & \NA & \NA \\
& Model-matched distance-weighted KNN & \NA & \NA & \NA & \NA & \NA \\
\bottomrule
\end{tabular}
\end{adjustbox}
\TableNote{Pending: dataset and split manifests are not frozen. Calibration Questions will report unique calibration IDs, not their three A/B/C inference records; the shared held-out-test $N$ must be stated once finalized. All rows within an evaluation-LLM block must use identical held-out IDs, with zero calibration/validation/test overlap. Cross-LLM rows are transfer diagnostics; the model-matched row is primary.}
\end{table}
